\documentclass[UTF8, 12pt, fontset=fandol, oneside]{ctexart}
\usepackage{researchnotes}

\title{各种平均数的概念与关系}
\author{}
\date{}

\begin{document}

\maketitle

平均数是把一组数压缩成一个代表值的方法. 在不同问题中，``代表''的含义并不相同：有时我们关心总量的均分，有时关心比例的连续增长，有时关心速度、效率等倒数型量，有时又希望较大的观测值得到更高的惩罚. 因此，平均数不是唯一的概念，而是一族根据问题结构选择的数值概括方式.

设
\[
  x_1,x_2,\cdots,x_n
\]
是一组实数. 若需要使用几何平均、调和平均等涉及乘积或倒数的平均数，通常还要求 $x_i>0$. 若带有权重，则设
\[
  w_i\ge0,\qquad \sum_{i=1}^n w_i=1.
\]
权重 $w_i$ 表示第 $i$ 个数据在平均过程中的相对重要性.

\section*{一、算术平均}

\noindent\textbf{定义}\quad 对实数 $x_1,x_2,\cdots,x_n$，其算术平均数定义为
\[
  A=\frac{x_1+x_2+\cdots+x_n}{n}.
\]
若带权重 $w_1,w_2,\cdots,w_n$，则加权算术平均为
\[
  A_w=\sum_{i=1}^n w_i x_i.
\]

算术平均的核心思想是\textbf{总量均分}. 若 $x_i$ 表示每个人拥有的物品数量、每次测量的误差、若干门课程的成绩等，把总和平均分给 $n$ 个对象，所得数值就是算术平均.

\noindent\textbf{基本性质}\quad 算术平均具有线性性质：
\[
  \frac{(ax_1+b)+\cdots+(ax_n+b)}{n}
  =aA+b.
\]
这说明对数据整体进行线性变换后，平均数也作同样的线性变换. 算术平均还满足
\[
  \sum_{i=1}^n (x_i-A)=0,
\]
即各数据相对平均数的偏差和为零.

\noindent\textbf{例}\quad 若三次测量结果为 $2,4,9$，则
\[
  A=\frac{2+4+9}{3}=5.
\]
若三门课学分分别为 $2,3,5$，成绩分别为 $80,90,86$，则加权平均成绩为
\[
  \frac{2}{10}\cdot80+\frac{3}{10}\cdot90+\frac{5}{10}\cdot86=86.
\]

\section*{二、几何平均}

\noindent\textbf{定义}\quad 对正数 $x_1,x_2,\cdots,x_n$，其几何平均数定义为
\[
  G=\sqrt[n]{x_1x_2\cdots x_n}.
\]
带权几何平均为
\[
  G_w=\prod_{i=1}^n x_i^{w_i}.
\]

几何平均的核心思想是\textbf{乘法结构的均分}. 若若干变化率连续相乘，或者若干比例共同决定最终结果，那么几何平均比算术平均更自然. 它回答的问题是：若每一步都用同一个倍率 $G$ 来替代原来的倍率 $x_1,\cdots,x_n$，怎样才能保持总乘积不变？

\noindent\textbf{对数视角}\quad 因为
\[
  \ln G=\frac{1}{n}\sum_{i=1}^n \ln x_i,
\]
所以几何平均就是先对数据取对数，再求算术平均，最后指数还原. 这也说明几何平均适合处理比例、倍数、增长率等量.

\noindent\textbf{例}\quad 某投资第一年增长 $20\%$，第二年下降 $10\%$，第三年增长 $50\%$. 三年的增长倍率分别为
\[
  1.2,\quad 0.9,\quad 1.5.
\]
等效的年均增长倍率为
\[
  G=\sqrt[3]{1.2\cdot0.9\cdot1.5}=\sqrt[3]{1.62}.
\]
因此等效年均增长率为
\[
  \sqrt[3]{1.62}-1.
\]
这里不能直接取 $20\%,-10\%,50\%$ 的算术平均，因为本金在每一年都会被前一年的变化重新缩放.

\section*{三、调和平均}

\noindent\textbf{定义}\quad 对正数 $x_1,x_2,\cdots,x_n$，其调和平均数定义为
\[
  H=\frac{n}{\frac{1}{x_1}+\frac{1}{x_2}+\cdots+\frac{1}{x_n}}.
\]
带权调和平均为
\[
  H_w=\frac{1}{\sum_{i=1}^n \frac{w_i}{x_i}}.
\]

调和平均的核心思想是\textbf{倒数结构的均分}. 当数据本身是``单位资源产生多少结果''或``单位结果消耗多少资源''这类比率时，若固定的是分子或分母中的某一总量，调和平均往往是合适的平均方式.

\noindent\textbf{例}\quad 一辆车去程速度为 $60$ 千米/小时，返程速度为 $40$ 千米/小时，且去程和返程路程相同. 平均速度不是
\[
  \frac{60+40}{2}=50,
\]
而是
\[
  H=\frac{2}{\frac{1}{60}+\frac{1}{40}}=48.
\]
这是因为平均速度等于总路程除以总时间，而时间与速度的倒数成正比.

\section*{四、平方平均}

\noindent\textbf{定义}\quad 对实数 $x_1,x_2,\cdots,x_n$，其平方平均数，也称均方根，定义为
\[
  Q=\sqrt{\frac{x_1^2+x_2^2+\cdots+x_n^2}{n}}.
\]
带权平方平均为
\[
  Q_w=\sqrt{\sum_{i=1}^n w_i x_i^2}.
\]

平方平均先把数据平方，再求算术平均，最后开平方. 由于平方会放大较大数的影响，所以平方平均对大值更敏感. 它常用于交流电有效值、误差大小、向量长度以及统计中的标准差相关计算.

\noindent\textbf{例}\quad 若误差为 $-1,2,-2$，其算术平均为
\[
  \frac{-1+2-2}{3}=-\frac{1}{3},
\]
这个数会因正负抵消而显得偏小. 若关心误差的典型大小，则可用平方平均：
\[
  Q=\sqrt{\frac{(-1)^2+2^2+(-2)^2}{3}}=\sqrt{3}.
\]

\section*{五、幂平均}

\noindent\textbf{定义}\quad 对正数 $x_1,x_2,\cdots,x_n$ 和实数 $p$，$p$ 阶幂平均定义为
\[
  M_p=
  \begin{cases}
    \left(\dfrac{x_1^p+x_2^p+\cdots+x_n^p}{n}\right)^{1/p}, & p\ne0,\\[1.2em]
    \sqrt[n]{x_1x_2\cdots x_n}, & p=0.
  \end{cases}
\]
带权形式为
\[
  M_p(w)=
  \begin{cases}
    \left(\displaystyle\sum_{i=1}^n w_i x_i^p\right)^{1/p}, & p\ne0,\\[1.2em]
    \displaystyle\prod_{i=1}^n x_i^{w_i}, & p=0.
  \end{cases}
\]

幂平均把许多常见平均数统一到同一个框架中：
\[
  M_1=A,\qquad M_0=G,\qquad M_{-1}=H,\qquad M_2=Q.
\]
当 $p$ 越大时，大的观测值影响越强；当 $p$ 越小时，小的观测值影响越强.

\noindent\textbf{幂平均不等式}\quad 若 $p<q$，则
\[
  M_p\le M_q.
\]
特别地，对正数 $x_1,\cdots,x_n$ 有
\[
  H\le G\le A\le Q.
\]
等号成立当且仅当
\[
  x_1=x_2=\cdots=x_n.
\]
这条链式不等式是比较各种平均数时最重要的基本事实之一.

\section*{六、最值平均与极限情形}

幂平均还包含两个极限情形：
\[
  \lim_{p\to+\infty}M_p=\max\{x_1,\cdots,x_n\},
\]
\[
  \lim_{p\to-\infty}M_p=\min\{x_1,\cdots,x_n\}.
\]
因此，最大值与最小值也可以看作广义平均数. 从这个角度看，幂平均构成一条从最小值逐渐过渡到最大值的连续族.

\noindent\textbf{直观解释}\quad 当 $p$ 很大时，$x_i^p$ 会极度放大最大的数，其他数对总和的贡献相对变小；当 $p$ 很小时，特别是趋向 $-\infty$ 时，最小的正数会通过负幂占据主导地位.

\section*{七、对数平均}

\noindent\textbf{定义}\quad 对两个不同正数 $a,b$，其对数平均定义为
\[
  L(a,b)=\frac{b-a}{\ln b-\ln a}.
\]
若 $a=b$，则定义
\[
  L(a,a)=a.
\]

对数平均自然出现在指数增长、热传导、平均温差等问题中. 它可以理解为函数 $f(x)=1/x$ 在区间 $[a,b]$ 上积分平均的倒数，因为
\[
  \frac{1}{b-a}\int_a^b \frac{1}{x}\,dx
  =\frac{\ln b-\ln a}{b-a},
\]
所以
\[
  L(a,b)=
  \left(
    \frac{1}{b-a}\int_a^b \frac{1}{x}\,dx
  \right)^{-1}.
\]

对两个正数，有如下典型关系：
\[
  G(a,b)\le L(a,b)\le A(a,b),
\]
其中
\[
  G(a,b)=\sqrt{ab},\qquad A(a,b)=\frac{a+b}{2}.
\]

\section*{八、积分平均}

\noindent\textbf{定义}\quad 若函数 $f$ 在区间 $[a,b]$ 上可积，则 $f$ 在 $[a,b]$ 上的平均值定义为
\[
  \overline{f}
  =\frac{1}{b-a}\int_a^b f(x)\,dx.
\]
它是连续对象上的算术平均. 离散情形中，我们把有限多个函数值相加再除以个数；连续情形中，求和由积分替代，个数由区间长度替代.

\noindent\textbf{平均值定理的联系}\quad 若 $f$ 在 $[a,b]$ 上连续，则存在 $\xi\in[a,b]$，使得
\[
  f(\xi)=\frac{1}{b-a}\int_a^b f(x)\,dx.
\]
这说明连续函数在区间上的平均值至少会被函数在某一点真正取到.

\section*{九、统计中的位置平均}

在统计中，平均数常用于描述数据的集中趋势. 但当数据含有异常值或明显偏斜时，算术平均可能并不稳健，此时还会使用中位数、众数、截尾平均等位置概括量.

\subsection*{1. 中位数}

\noindent\textbf{定义}\quad 将数据从小到大排列. 若样本量为奇数，则中位数为正中间的数；若样本量为偶数，则中位数通常取中间两个数的算术平均.

中位数的特点是只关心排序位置，不关心极端值有多大. 因此它比算术平均更抗异常值.

\noindent\textbf{例}\quad 数据
\[
  2,\quad 3,\quad 4,\quad 5,\quad 100
\]
的算术平均为 $22.8$，但中位数为 $4$. 若用一个数表示这组数据的典型水平，中位数可能比算术平均更合理.

\subsection*{2. 众数}

\noindent\textbf{定义}\quad 众数是数据中出现次数最多的值. 如果多个值出现次数并列最多，则可能有多个众数.

众数适合描述类别型数据或离散数据的最常见取值. 例如服装尺码、调查选项、考试分数段等问题中，众数常比算术平均更有解释意义.

\subsection*{3. 截尾平均}

\noindent\textbf{定义}\quad 截尾平均是先删去一部分最大值和最小值，再对剩余数据求算术平均. 例如 $10\%$ 截尾平均就是删去最小的 $10\%$ 和最大的 $10\%$ 数据后，对中间部分求平均.

截尾平均介于算术平均和中位数之间：它保留了大部分数值信息，同时减弱了极端值的影响.

\section*{十、移动平均}

\noindent\textbf{定义}\quad 对时间序列
\[
  x_1,x_2,\cdots,x_N,
\]
长度为 $k$ 的简单移动平均定义为
\[
  m_t=\frac{x_{t-k+1}+x_{t-k+2}+\cdots+x_t}{k},
  \qquad t=k,k+1,\cdots,N.
\]

移动平均用于平滑短期波动、显示长期趋势. 它并不是把整组数据压缩成一个数，而是用一个滑动窗口不断计算局部平均.

\noindent\textbf{指数移动平均}\quad 指数移动平均用递推式
\[
  s_t=\alpha x_t+(1-\alpha)s_{t-1},\qquad 0<\alpha<1,
\]
给近期数据更高权重. $\alpha$ 越大，平均值对新数据越敏感；$\alpha$ 越小，曲线越平滑.

\section*{十一、如何选择平均数}

选择平均数时，关键不是看公式是否熟悉，而是看问题中的量如何组合.

\noindent\textbf{总量可加时用算术平均}\quad 若总量由各项直接相加得到，例如人数、分数、长度、收入总额等，通常使用算术平均.

\noindent\textbf{倍率相乘时用几何平均}\quad 若总效果由若干倍率相乘得到，例如复利收益、连续增长率、相对风险倍数等，通常使用几何平均.

\noindent\textbf{倒数量相加时用调和平均}\quad 若问题中总时间、单位成本、单位产出等与数据的倒数成正比，常使用调和平均.

\noindent\textbf{关心大小而不关心符号时用平方平均}\quad 若正负抵消会掩盖问题，例如误差强度、波动幅度、信号有效值等，可考虑平方平均.

\noindent\textbf{异常值明显时考虑中位数或截尾平均}\quad 若少量极端观测值会严重拉动算术平均，则中位数和截尾平均通常更稳健.

\section*{十二、常见平均数对照表}

\begin{center}
\begin{tabularx}{\textwidth}{L{3.1cm}Y Y}
\toprule
名称 & 公式 & 适用直观 \\
\midrule
算术平均 & $\displaystyle A=\frac{1}{n}\sum_{i=1}^n x_i$ & 总量均分，适合可加量 \\
几何平均 & $\displaystyle G=\left(\prod_{i=1}^n x_i\right)^{1/n}$ & 乘法均分，适合倍率和比例 \\
调和平均 & $\displaystyle H=\frac{n}{\sum_{i=1}^n 1/x_i}$ & 倒数均分，适合速度、效率等比率 \\
平方平均 & $\displaystyle Q=\sqrt{\frac{1}{n}\sum_{i=1}^n x_i^2}$ & 放大大值，适合误差和有效值 \\
幂平均 & $\displaystyle M_p=\left(\frac{1}{n}\sum_{i=1}^n x_i^p\right)^{1/p}$ & 统一多种平均，$p$ 控制大小值影响 \\
对数平均 & $\displaystyle L(a,b)=\frac{b-a}{\ln b-\ln a}$ & 指数变化、积分平均相关问题 \\
积分平均 & $\displaystyle \overline{f}=\frac{1}{b-a}\int_a^b f(x)\,dx$ & 连续对象上的平均值 \\
中位数 & 排序后位于中间位置的值 & 稳健描述集中位置 \\
截尾平均 & 删除两端极值后求算术平均 & 兼顾数值信息和稳健性 \\
\bottomrule
\end{tabularx}
\end{center}

\section*{十三、一个统一观点}

许多平均数都可以看作先把数据通过某个函数变换，再求算术平均，最后用反函数变回原尺度. 若 $\varphi$ 是严格单调函数，则可定义
\[
  M_{\varphi}
  =\varphi^{-1}
  \left(
    \frac{1}{n}\sum_{i=1}^n \varphi(x_i)
  \right).
\]
这称为拟算术平均. 取不同的 $\varphi$，可得到不同平均数：
\[
  \varphi(x)=x \quad\Longrightarrow\quad M_{\varphi}=A,
\]
\[
  \varphi(x)=\ln x \quad\Longrightarrow\quad M_{\varphi}=G,
\]
\[
  \varphi(x)=\frac{1}{x} \quad\Longrightarrow\quad M_{\varphi}=H,
\]
\[
  \varphi(x)=x^2 \quad\Longrightarrow\quad M_{\varphi}=Q
  \quad (x_i\ge0).
\]

这个观点说明：平均数的本质不是简单地``相加再除以个数''，而是先决定问题中真正线性的尺度是什么. 在普通数量尺度上线性，就得到算术平均；在对数尺度上线性，就得到几何平均；在倒数尺度上线性，就得到调和平均；在平方尺度上线性，就得到平方平均.

\section*{十四、小结}

平均数的选择应服从问题结构. 算术平均适合可加总量，几何平均适合乘法倍率，调和平均适合倒数型比率，平方平均适合衡量强度和波动，幂平均提供统一的比较框架，对数平均和积分平均则常出现在连续变化与分析问题中. 在统计应用中，中位数、众数和截尾平均提醒我们：一个代表值不仅要有公式上的合理性，也要对数据分布和实际语境负责.

\end{document}
